\documentclass[12pt]{amsart}
%prepared in AMSLaTeX, under LaTeX2e
\addtolength{\oddsidemargin}{-1.1in} 
\addtolength{\evensidemargin}{-1.1in}
\addtolength{\topmargin}{-.9in}
\addtolength{\textwidth}{1.75in}
\addtolength{\textheight}{1.75in}

\renewcommand{\baselinestretch}{1.05}

\usepackage{verbatim} % for "comment" environment

\usepackage{palatino}

\usepackage[final]{graphicx}

\usepackage{tikz}
\usetikzlibrary{positioning}

\usepackage{bm,enumitem,xspace,fancyvrb}

\newtheorem*{thm}{Theorem}
\newtheorem*{defn}{Definition}
\newtheorem*{example}{Example}
\newtheorem*{problem}{Problem}
\newtheorem*{remark}{Remark}

\DefineVerbatimEnvironment{mVerb}{Verbatim}{numbersep=2mm,frame=lines,framerule=0.1mm,framesep=2mm,xleftmargin=4mm,fontsize=\footnotesize}

% macros
\usepackage{amssymb}
\newcommand{\bA}{\mathbf{A}}
\newcommand{\bB}{\mathbf{B}}
\newcommand{\bE}{\mathbf{E}}
\newcommand{\bF}{\mathbf{F}}
\newcommand{\bJ}{\mathbf{J}}

\newcommand{\bb}{\mathbf{b}}
\newcommand{\bc}{\mathbf{c}}
\newcommand{\bd}{\mathbf{d}}
\newcommand{\bi}{\mathbf{i}}
\newcommand{\bj}{\mathbf{j}}
\newcommand{\bk}{\mathbf{k}}
\newcommand{\br}{\mathbf{r}}
\newcommand{\bv}{\mathbf{v}}
\newcommand{\bw}{\mathbf{w}}
\newcommand{\bx}{\mathbf{x}}

\newcommand{\bzero}{\bm{0}}

\newcommand{\eps}{\epsilon}
\newcommand{\grad}{\nabla}
\newcommand{\ip}[2]{\ensuremath{\left<#1,#2\right>}}
\newcommand{\lam}{\lambda}
\newcommand{\lap}{\triangle}

\newcommand{\Null}{\operatorname{null}}
\newcommand{\rank}{\operatorname{rank}}
\newcommand{\range}{\operatorname{range}}
\newcommand{\trace}{\operatorname{tr}}

\newcommand{\RR}{\mathbf{R}}

\newcommand{\ZZ}{\mathbb{Z}}

\newcommand{\prob}[1]{\bigskip\noindent\textbf{#1.}\quad }
\newcommand{\exer}[2]{\prob{Exercise #2 on page #1}}

\newcommand{\pts}[1]{(\emph{#1 pts}) }
\newcommand{\epart}[1]{\bigskip\noindent\textbf{(#1)}\quad }
\newcommand{\ppart}[1]{\,\textbf{(#1)}\quad }

\newcommand{\Matlab}{\textsc{Matlab}\xspace}
\newcommand{\ds}{\displaystyle}

\def\vectwo#1#2{\begin{bmatrix}#1\\#2\end{bmatrix}}
\def\vecthree#1#2#3{\begin{bmatrix}#1\\#2\\#3\end{bmatrix}}
\def\vecfour#1#2#3#4{\begin{bmatrix}#1\\#2\\#3\\#4\end{bmatrix}}
\def\bbm{\begin{bmatrix}}
\def\ebm{\end{bmatrix}}


\begin{document}
\scriptsize \noindent Math 314 Linear Algebra \hfill Section 3.4 and 3.5 Examples 
\normalsize\medskip


\thispagestyle{empty}
\begin{enumerate}
\item For each matrix $A$ below, its columns will be labelled $\mathbf{a}_i$.\\
	\begin{itemize}
	\item Use $R$ to write the columns of $A$ as linear combinations of each other or state that this is not possible.
	\item Find a \emph{basis} for $C(A)$, the column space of $A$ and  $N(A),$ the null space of A. 
	\item Determine the \emph{dimensions} of $C(A)$ and  $N(A)$.
	\end{itemize}
	\begin{enumerate}
	\item $A \quad \to \quad R = \begin{bmatrix} 1 & 0 & 2 \\ 0 & 1 & 3 \\ 0 & 0 & 0 \end{bmatrix}$
	\vfill
	\item $A \quad \to \quad R = \begin{bmatrix} 1 & 0 & 2 &0&0\\ 0 & 1 & 0&0&3 \\ 0 & 0 & 0&1&-4 \\0&0&0&0&0\end{bmatrix}$
	\vfill
	\item $A \quad \to \quad R = \begin{bmatrix} 1 & 0 & 0 \\ 0 & 1 & 0 \\ 0 & 0 & 1 \end{bmatrix}$
	\vspace{.5in}
	\end{enumerate}
\newpage
\item (summary of previous page) Suppose $A$ is an $m \times n$ matrix of rank $r$.\\
	\begin{enumerate}
	\item The dimension of $C(A)$ is $\underline{\qquad\qquad}$ and a basis for $C(A)$ can be found by\\
	\medskip
	\medskip
	\item The dimension of $N(A)$ is $\underline{\qquad\qquad}$ and a basis for $N(A)$ can be found by\\
	\medskip
	\medskip
	\end{enumerate}
\item Find a basis for the \emph{row space of $A$} where $A = \begin{bmatrix} 1 & 0 & 2 \\ 1 & 1 & 5 \\ 2 & -2 & -2 \end{bmatrix}.$
\vfill
Revisit 1b.
\end{enumerate}
\end{document}

\item $T_2=\left\{ \mathbf{v}_1= \vecthree{12}{-9}{1},\mathbf{v}_2= \vecthree{-10}{-1}{3},\mathbf{v}_3= \vecthree{5}{-5}{12} \right\}$\\

\medskip
\noindent  $A = \begin{bmatrix} 12 & -10 & 5 \\ -9 & -1 & -5 \\ 1 & 3 & 12 \end{bmatrix} \qquad \to \qquad R = \begin{bmatrix} 1 & 0 & 0 \\ 0 & 1 & 0 \\ 0 & 0 & 1 \end{bmatrix}$
\vfill
\item $T_3=\left\{ \mathbf{v}_1= \vecthree{12}{-9}{1},\mathbf{v}_2= \vecthree{-10}{-1}{3},\mathbf{v}_3= \vecthree{5}{-5}{12}, \mathbf{v}_4= \vecthree{7}{-15}{16} \right\}$\\
\vfill


\newpage

\item $T_4=\left\{ \mathbf{v}_1= \vectwo{2}{2},\mathbf{v}_2= \vectwo{-1}{1},\mathbf{v}_3= \vectwo{5}{-1},\mathbf{v}_3=\vectwo{2}{6} \right\}$\\

\medskip
\noindent $A = \begin{bmatrix} 2 & -1 & 5 & 2 \\ 2 & 1 & -1 & 6 \end{bmatrix} \qquad \to \qquad R = \begin{bmatrix} 1 & 0 & 1 & 2 \\ 0 & 1 & -3 & 2 \end{bmatrix}$
\vfill

\item $T_5=\left\{ \mathbf{v}_1= \vecfour{2}{-1}{5}{2},\mathbf{v}_2= \vecfour{2}{1}{-1}{6} \right\}$\\

\medskip
\noindent $A = \begin{bmatrix} 2 & 2 \\ -1 & 1 \\ 5 & -1 \\ 2 & 6 \end{bmatrix} \qquad \to \qquad R = \begin{bmatrix} 1 & 0 \\ 0 & 1 \\ 0 & 0 \\ 0 & 0 \end{bmatrix}$
\vfill

\end{enumerate}
\end{document}
